% part: intuitionistic-logic
% chapter: propositions-as-types
% section: introduction

\documentclass[../../../include/open-logic-section]{subfiles}

\begin{document}

\olfileid[id]{int}{pty}{int}

\olsection{Pendahuluan}

Kalkulus lambda adalah sebuah model komputasi sederhana. Kalkulus ini
bekerja pada term yang dibangun dari variabel. Jika $N$ dan $M$ adalah
term, maka $(NM)$ juga merupakan term (``$N$ diterapkan kepada~$M$'').
Jika $N$ merupakan term, maka $\lambd[x][N]$ merupakan term. Jika $N$
memuat variabel~$x$, kemunculan-kemunculan $x$ yang bersesuaian dalam
$\lambd[x][N]$ diikat oleh operator $\lambd[x]$, sehingga tidak lagi
bebas. Term berbentuk $(\lambd[x][N]M)$ dikatakan
berkontraksi-$\beta$ menjadi $\Subst{N}{M}{x}$ (hasil penggantian setiap
kemunculan bebas $x$ dalam~$N$ dengan $M$). Term $N$
berkonversi-$\beta$ menjadi~$N'$ jika $N'$ dihasilkan dari $N$ dengan
mengontraksi-$\beta$ suatu subterm, dan kita menulis $N \redone N'$.
Komputasi dalam kalkulus-$\lambd$ adalah suatu barisan $N \redone N_1
\redone N_2 \redone \dots$. Jika barisan ini berakhir pada term~$N_k$
yang tidak memiliki subterm yang dapat dikontraksi-$\beta$, term itu
disebut \emph{normal} atau bentuk normal dari~$N$. Bayangkan komputasi
dalam kalkulus-$\lambd$ sebagai barisan semacam itu. Komputasi berhenti
jika menghasilkan bentuk normal, dan bentuk normal tersebut merupakan
hasil komputasinya. Model komputasi sederhana ini ternyata
lengkap-Turing: setiap fungsi rekursif parsial dapat dihitung oleh suatu
term dalam kalkulus-$\lambd$.

Haskell Curry dan William Howard secara terpisah menemukan kemiripan erat
antara deduksi alami dan kalkulus-$\lambd$. Secara intuitif, term
$\lambd[x][N]$ adalah fungsi: $x$ merupakan argumennya; $N$ memberi tahu
kita cara menghitung nilai ketika $x$ diberikan. Term
$(\lambd[x][N]M)$ kemudian merupakan penerapan fungsi $\lambd[x][N]$
kepada argumen~$M$, dan kontraksi-$\beta$ memberi tahu kita cara
mengevaluasinya: gantikan $x$ dalam $N$ dengan~$M$ (lalu terus evaluasi
term yang dihasilkan). Sekarang bayangkan semua itu sebagai fungsi dengan
domain dan daerah nilai yang tetap. Jika domain $\lambd[x][N]$ adalah
$A$ dan daerah nilainya~$B$, variabel~$x$ seharusnya dipandang hanya
menerima argumen dari~$A$, sedangkan $N$ menghasilkan nilai dari~$B$.
Karena itu, $(\lambd[x][N]M)$ seharusnya bermakna hanya jika
$\lambd[x][N]$ merupakan fungsi $A \to B$ dan $M \in A$. Jika informasi
ini ditambahkan kepada kalkulus-$\lambd$, kita memperoleh
kalkulus-$\lambd$ \emph{bertipe}. Dalam kalkulus-$\lambd$ bertipe, tidak
setiap ekspresi $(\lambd[x][N]M)$ sah, melainkan hanya ekspresi yang
``tipe'' dari $\lambd[x][N]$ dan dari $M$ saling cocok, yakni
$\lambd[x][N]$ bertipe $A \lif B$ dan $M$ bertipe~$A$. Tipe
$\lambd[x][N]$ ialah $A \to B$ hanya jika tipe $x$ ialah $A$ dan tipe
$N$ ialah~$B$. Secara umum, tidak setiap ekspresi $(M_1M_2)$ sah. Hanya
term $(M_1M_2)$ dengan $M_1$ bertipe $A \to B$ dan $M_2$ bertipe~$A$
yang sah; jika demikian, $(M_1M_2)$ bertipe~$B$.

Aturan pengetikan ini jelas bersesuaian dengan !!{derivation}s dalam
deduksi alami, dengan tipe direpresentasikan oleh !!{formula}s dan
!!{derivation}s itu sendiri oleh term-$\lambd$. Misalnya,
!!a{derivation} atas $!A \lif !B$ bersesuaian dengan term
$\lambd[\typeof{x}{!A}][N]$, yang mempunyai tipe fungsi
$!A \lif !B$. Aturan inferensi deduksi alami bersesuaian dengan aturan
pengetikan dalam kalkulus lambda bertipe, misalnya,
\begin{prooftree}
  \AxiomC{$\Discharge{!A}{x}$}
  \DeduceC{$!B$}
  \DischargeRule{\Intro\lif}{x}
  \UnaryInfC{$!A \lif !B$}
  \DisplayProof\bottomAlignProof
  \quad\text{bersesuaian dengan}\quad
  \Axiom$x: !A \fCenter N: !B$
  \UnaryInf$ \fCenter \lambd[\typeof{x}{!A}][N] : !A \lif !B$
\end{prooftree}
dengan aturan di kanan berarti bahwa jika $x$ bertipe~$!A$ dan $N$
bertipe~$!B$, maka $\lambd[\typeof{x}{!A}][N]$ bertipe~$!A \lif !B$.

Aturan $\Elim\lif$ bersesuaian dengan aturan pengetikan bagi term
aplikasi, yakni,
\[
  \AxiomC{$!A \lif !B$}
  \AxiomC{$!A$}
  \RightLabel{\Elim\lif}
  \BinaryInfC{$!B$}
  \DisplayProof
  \quad\text{bersesuaian dengan}\quad
  \Axiom$\fCenter M_1: !A \lif !B$
  \Axiom$\fCenter M_2: !A$
  \BinaryInf$ \fCenter (M_1M_2) : !B$
  \DisplayProof
\]
Jika aturan \Intro\lif{} langsung diikuti oleh aturan \Elim\lif{},
!!{derivation} tersebut dapat disederhanakan:
\begin{gather*}
  \AxiomC{$\Discharge{!A}{x}$}
  \DeduceC{$!B$}
  \DischargeRule{\Intro\lif}{x}
  \UnaryInfC{$!A \lif !B$}
  \AxiomC{}
  \DeduceC{$!A$}
  \RightLabel{\Elim\lif}
  \BinaryInfC{$!B$}
  \DisplayProof
  \quad\redone\quad
  \AxiomC{}
  \DeduceC{$!A$}
  \DeduceC{$!B$}
  \DisplayProof
\end{gather*}
yang bersesuaian dengan kontraksi-$\beta$ pada term lambda
\[
(\lambd[\typeof{x}{!A}][N]M) \quad\redone\quad
\Subst{N}{M}{x}.
\]

Korespondensi serupa berlaku antara aturan untuk $\land$ dan tipe
``produk,'' serta antara aturan untuk $\lor$ dan tipe ``jumlah.''

Korespondensi antara term dalam kalkulus lambda bertipe sederhana dan
!!{derivation}s deduksi alami ini disebut korespondensi ``Curry--Howard,''
``bukti sebagai program,'' atau ``proposisi sebagai tipe.'' Selain
!!{formula}s (proposisi) yang bersesuaian dengan tipe, dan bukti yang
bersesuaian dengan term, kita dapat merangkum korespondensi tersebut
sebagai berikut:
\begin{center}
  \begin{tabular}{c c}
    logika & program \\
    \hline
    proposisi/!!{formula} & tipe \\
    bukti & term \\
    asumsi & variabel \\
    asumsi yang dikosongkan & variabel terikat \\
    asumsi yang tidak dikosongkan & variabel bebas \\
    $!A \lif !B$ & tipe fungsi \\
    $!A \land !B$ & tipe produk \\
    $!A \lor !B$ & tipe jumlah \\
    $\lfalse$ & tipe kosong \\
    \hline
  \end{tabular}
\end{center}

Korespondensi Curry--Howard merupakan salah satu landasan asisten bukti
otomatis dan pemeriksa tipe bagi program, sebab memeriksa bukti yang
menyaksikan suatu proposisi (seperti yang kita lakukan di atas) sama
dengan memeriksa apakah suatu program (term) mempunyai tipe yang
dideklarasikan.
\end{document}
